
\chapter{Embedding Classical Field Theories into the Lamphrodynamic Plenum}
\label{ch:embedding}

\subsection{Introduction}

Classical field theory is usually formulated in terms of sections of bundles over a fixed spacetime manifold, together with variational principles determining the Euler--Lagrange equations. In the present volume we have recast field theory as arising from derived moduli spaces equipped with a $(-1)$--shifted symplectic structure. The goal of this chapter is to relate the ordinary formulation to the lamphrodynamic formalism by constructing a functor from a suitable category of classical field theories to the category representing RSVP field models.

We adopt the interpretative stance that classical theories embed \emph{functorially} into RSVP. That is, there exists an $(\infty,1)$--functor
[
F \colon \CFTh \longrightarrow \RSVP,
]
which is fully faithful (or at least conservative), and whose essential image contains precisely those lamphrodynamic models that arise from standard classical physics. In particular, the familiar theories of scalar fields, gauge fields, and Einstein gravity all appear as objects inside RSVP rather than as limits or degenerations.

\begin{remark}
This point of view corresponds to Option~C of the discussion ending the previous chapter: RSVP is a categorical enlargement of classical theory, not merely a generalization nor an approximation scheme.
\end{remark}

\subsection{The category of classical field theories}

Let $X$ be a smooth oriented spacetime manifold and let $E \to X$ be a smooth vector bundle (or principal bundle). A classical field configuration is a smooth section $\phi \colon X \to E$. A classical Lagrangian density $\mathcal{L}$ is a smooth function of the jet bundle $J^\infty(E)$, and the classical action is
[
S_{\mathrm{class}}[\phi] = \int_X \mathcal{L}(\phi, \nabla\phi, \ldots ), .
]

\begin{definition}
Define the $(\infty,1)$--category $\CFTh$ whose objects are triples $(X,E,\mathcal{L})$ and whose morphisms are field--bundle--variational morphisms preserving the Lagrangian structure.
\end{definition}

This category includes scalar field theory, Yang--Mills theory, and Einstein gravity, among others. It constitutes the standard domain of classical theoretical physics.

\subsection{Construction of the embedding functor}

We now define a functor
[
F \colon \CFTh \longrightarrow \RSVP.
]
The idea is that a classical field $\phi \colon X \to E$ is assigned a lamphrodynamic triple $(\Phi,\mathbf v,S)$ via a canonical lift. The potential $\Phi$ encodes the classical field value, while $\mathbf v$ and $S$ are obtained from trivial (or canonical) background data determined by the geometry of $X$ and the structure of $E$. Formally, set
[
F(X,E,\mathcal{L}) = \Map(X,\Plenum_{\Phi,\mathbf v,S}),
]
where $\Plenum_{\Phi,\mathbf v,S}$ is the lamphrodynamic moduli stack introduced previously, equipped with the $(-1)$--shifted symplectic structure constructed in Chapter~6.

\begin{proposition}
The assignment $F$ extends to an $(\infty,1)$--functor $F \colon \CFTh \to \RSVP$, preserving homotopy equivalences and sending classical Lagrangian morphisms to lamphrodynamic morphisms.
\end{proposition}

\begin{proof}[Sketch of proof]
A morphism of classical field theories induces a morphism of derived mapping stacks. The shifted symplectic forms are natural with respect to these maps by functoriality of the AKSZ construction. Hence the construction is functorial. The full proof uses the theory of tangent complexes discussed in Chapter~2.
\end{proof}

\subsection{Examples}

\subsubsection{Scalar fields}

Let $(X,E,\mathcal{L})$ be a scalar field theory with $E = X \times \mathbb{R}$ and
[
\mathcal{L} = \frac12(\partial\phi)^2 - V(\phi).
]
Then
[
F(X,E,\mathcal{L}) = \Map(X,\Plenum),
]
with $F(\phi) = (\Phi,\mathbf v,S)$ given by $\Phi(x) = \phi(x)$ and $\mathbf v,S$ trivial. The Euler--Lagrange equations become the $\Phi$--component of the lamphrodynamic equations.

\subsubsection{Yang--Mills fields}

For Yang--Mills theory the classical gauge potential determines a component of $\Phi$ via a fixed embedding $\mathfrak{g}\hookrightarrow\text{scalar sector}$, while $\mathbf v$ and $S$ remain trivial. The shifted symplectic form lifts from the moduli of gauge fields by AKSZ.

\subsubsection{Einstein gravity}

In the case of Einstein gravity, the Lorentzian metric determines $\Phi$ through a chosen geometric invariant such as curvature, and the components $\mathbf v$ and $S$ encode kinematic and entropic contributions from spacetime geometry. Details require the constructions of Chapter~12.

\subsection{Conclusion}

The functor $F$ embeds classical field theory into the RSVP formalism as a distinguished substructure. In this manner, the classical paradigm appears as a special case of a more general lamphrodynamic framework, and quantization can proceed categorically at the level of RSVP rather than theory by theory.

\begin{remark}
This approach allows one to regard classical field theories and RSVP models within a unified mathematical language, and it clarifies the relationship between classical physics and the derived geometric structures developed in the preceding chapters.
\end{remark}

